\begin{longtable}{p{3cm} p{3cm} p{10cm}}
    \caption{Campos de un mensaje Syslog según RFC 3164}
    \label{tab:campos-rfc3164} \\

    \toprule
    \textbf{Campo} & \textbf{Valor} & \textbf{Descripción} \\
    \midrule
    \endfirsthead

    \multicolumn{3}{c}{\textit{Tabla \thetable. Continuación}} \\
    \toprule
    \textbf{Campo} & \textbf{Valor} & \textbf{Descripción} \\
    \midrule
    \endhead

    \midrule
    \multicolumn{3}{r}{\textit{Continúa en la siguiente página}} \\
    \endfoot

    \bottomrule
    \endlastfoot

    \texttt{PRI} &
    \texttt{<34>} &
    Indica la prioridad del mensaje y está compuesto por los valores de Facility y Severity. \\

    \texttt{TIMESTAMP} &
    \texttt{Oct 11 22:14:15} &
    Indica la fecha y hora en que se generó el evento. Utiliza el formato de mes, día, hora, minutos y segundos. \\

    \texttt{HOSTNAME} &
    \texttt{myhost} &
    Identifica el sistema o dispositivo que originó el mensaje. \\

    \texttt{TAG} &
    \texttt{su} &
    Permite identificar la aplicación o proceso asociado con la generación del mensaje. \\

    \texttt{PID} &
    \texttt{[1234]} &
    Identifica el proceso que generó el evento cuando el sistema emisor proporciona esta información. \\

    \texttt{MSG} &
    \texttt{'su root' failed for jdoe on /dev/pts/2} &
    Contiene la información específica del evento generado por la aplicación o dispositivo. \\

\end{longtable}